\chapter{The Residue Refuses To Disappear}

\section{The antagonist and the treasure}

The last chapter closed a loop and found the count unchanged. But it ended on a warning: some loops come back
\emph{altered}, and the alteration is the whole point. This chapter meets that alteration and names it. The residue is the
part of a round trip that fails to cancel --- the leftover a closed circuit deposits that no relabeling can wash out. It is
at once the antagonist of the story, because it will not go away, and its treasure, because everything the instrument later
measures is read off it. On this face the residue has a precise name: it is a curvature, and a curvature that is closed but
not exact.

The geometry is clean. Carry a vector around a small closed loop by parallel transport under a connection $\nabla$, and ask
whether it returns to itself. In flat space it does; in general it does not, and the discrepancy is the curvature,
\[
  [\nabla_\mu, \nabla_\nu]\,V^\rho \;=\; R^{\rho}{}_{\sigma\mu\nu}\,V^\sigma .
\]
The commutator of two transports is the failure of the loop to close, and $R^{\rho}{}_{\sigma\mu\nu}$ is exactly that failure
made into a tensor. The same shape governs the gauge field: the field strength is the curvature of the connection one-form
$A$,
\[
  F \;=\; \mathrm{d}A, \qquad F_{\mu\nu} = \partial_\mu A_\nu - \partial_\nu A_\mu,
\]
and $F$ measures the phase a state picks up on a closed circuit. In both cases the residue is what a loop deposits: transport
around and back, and the difference from the identity is a curvature that will not vanish.

What makes the residue an \emph{antagonist} --- a leftover that refuses to be discarded --- is stated exactly in the language
of forms. The field strength is closed, $\mathrm{d}F = 0$, which is the Bianchi identity; the question is whether it is also
exact, whether $F = \mathrm{d}A$ globally for a single-valued $A$. When it is not, the residue is cohomologically nontrivial:
its class $[F]$ in de Rham cohomology is a nonzero element, and no choice of gauge, no smooth relabeling of the potential,
can set it to zero. A closed form that is not exact is precisely a residue that cannot be discarded --- locally it looks like
pure gauge, a thing you could subtract away, but globally the loop integral $\oint A$ detects it and the alteration survives.
The holonomy is the residue you can measure by going around.

There is a physical reading laid over this geometry, and it is marked as a reading. Aharonov and Bohm showed that a charged
state carried around a region of nonzero flux returns with a phase $\exp\!\big(\tfrac{ie}{\hbar}\oint A\big)$ shifted by the
enclosed residue, even where the field strength itself vanishes along the path~\cite{aharonov1959}. That the loop's
alteration is \emph{physically observable} --- an interference fringe that shifts, a thing a bench reads --- is the physics
gauge's confirmation that the residue is not a bookkeeping artifact but the world's own leftover. The construction supplies
the closed-not-exact curvature; the identification of its loop integral with an observable phase is the reading, and the
laboratory has cashed it.

In the two verbs, the residue is what the funge cannot fully bag. The loop tries to gather its start and its end into one
class --- to funge the round trip into ``no change'' --- and cannot, because a covariant curvature stands between them that no
tange of the labels can reach. The residue is the world's, funged and invariant; our relabelings, the contravariant tange,
slide over it without touching it. That is why it is both antagonist and treasure: it is the one thing the relabeling cannot
alter, which makes it the one thing worth measuring. The count of the last chapter was the loop that closed; the residue is
the loop that does not, and the whole instrument from here is the reading of that difference.

\section{The lie of discarding}

The residue is a leftover, and leftovers invite disposal. The most natural error an instrument can make --- and the one the
physics of the last century spent enormous effort learning not to make --- is to throw the residue away: to drop the term
that will not cancel, subtract the piece that will not vanish, and present the tidied remainder as the answer. This section
is about why that disposal is a lie, and why the instrument's honesty consists precisely in refusing it. The residue is not
noise to be cleaned off a signal. It \emph{is} the signal, and discarding it discards the measurement.

The temptation has a standard geometric form: the boundary term. Integrate a total derivative over a region and it collapses
to the boundary, $\int_\Omega \mathrm{d}\omega = \oint_{\partial\Omega} \omega$, by Stokes' theorem; and it is tempting, when
a boundary term is inconvenient, to declare the fields to fall off fast enough that $\oint_{\partial\Omega}\omega = 0$ and
drop it. Sometimes that is legitimate. But when the residue is topological --- when $\oint A$ is the holonomy of a
nontrivial class --- the boundary term is not zero and cannot be assumed away, and setting it to zero by decree is not an
approximation but a falsification. The lie of discarding is the assumption, made for convenience, that a leftover which
happens to be inconvenient is therefore negligible. The instrument's discipline is to carry every boundary term until the
geometry, not the convenience, decides whether it vanishes.

The same lie wears a subtler dress in the quantum field, where the residue often arrives as an infinity, and the crude
response is to subtract it and never speak of it again. But the honest accounting is that the divergent piece is a real
feature of the naive description, and what the residue does under an honest hand is get \emph{carried} --- tracked through
the calculation, absorbed into a redefinition that is itself observable, never simply deleted. The difference between a lie
and a measurement here is the difference between subtracting the residue and \emph{following} it. A subtracted residue
leaves no trace and answers no question; a followed residue reappears, later, as a physical shift the bench can read. The
whole method of the instrument is to be the second kind of device --- the one that carries the leftover rather than the one
that quietly disposes of it.

And here is the deepest reason the disposal is a lie: the residue is exactly the part that is \emph{invariant}, and to
discard an invariant is to discard the frame-independent content and keep only the labels. The curvature $R^\rho{}_{\sigma
\mu\nu}$ and the class $[F]$ are the covariant, funged quantities --- the world's, unchanged by any relabeling. The parts you
\emph{can} freely alter are the tange, the contravariant labels, the gauge; those carry no physics and disposing of them
costs nothing. But the residue is not among them. To throw the residue away is to throw away the one thing a change of frame
could not touch --- to keep the arbitrary and discard the real. That inversion is the precise anatomy of the crank's error,
and the instrument is built to make it structurally impossible: it cannot dispose of the residue without disposing of its
own reading.

In the two verbs, the honesty is a refusal to over-tange. To discard the residue is to select it away --- to tange the
inconvenient leftover out of the account by hand --- and the instrument declines that selection, because the residue is a
funge, the world's covariant leftover, and no honest tange reaches it. What the instrument may tange away is only the labels,
the gauge, the arbitrary contravariant frame; what it must carry is the invariant funge, the curvature that stays. So the
lie of discarding and the honesty of carrying are, exactly, the difference between tanging away the world's leftover and
tanging away only our own labels. The instrument keeps the residue because keeping it is the only way to keep the
measurement.

\section{What the residue will become}

The residue has been named as a curvature and defended against disposal. This closing section says what it is \emph{for} ---
what a carried residue becomes once the instrument reads it at a particular place. The answer is the plan of the whole book:
one leftover, read at different sites of the construction, becomes mass, becomes charge, becomes phase, becomes curvature
itself. The residue is not many things. It is one invariant, and the physical names are the readings it takes on where the
apparatus meets it.

Follow the one curvature into its readings. Contract it and it becomes the source of gravity: the Ricci tensor $R_{\mu\nu} =
R^\rho{}_{\mu\rho\nu}$, tied to matter by Einstein's equation $R_{\mu\nu} - \tfrac{1}{2}R\,g_{\mu\nu} = 8\pi\,T_{\mu\nu}$,
reads the residue as \emph{mass}-energy~\cite{einstein1915,einstein1916}. Take the gauge curvature's flux and it becomes
\emph{charge}: the integral of $F$ over a closed surface counts the source it encloses, $\oint \star F \sim e$. Take its loop
integral and it becomes \emph{phase}: $\exp\!\big(\tfrac{ie}{\hbar}\oint A\big)$, the holonomy of the last section. Leave it
uncontracted, as the full $R^\rho{}_{\sigma\mu\nu}$, and it is \emph{curvature} in the plainest sense, the tidal bending of
the geometry. Four physical names, and beneath them a single object read at four sites --- the same residue the loop failed
to shed.

It must be marked exactly which of this is built and which is read. That the construction carries an invariant residue, a
closed-not-exact curvature that survives every loop --- that is the construction's, derived. That this residue \emph{is}
mass here, charge there, phase elsewhere --- that is the physical gauge's reading laid over the invariant, and it is marked
as a reading. The instrument does not prove that its leftover is the electron's charge; it proves that it carries an
invariant leftover, and the physics recognizes, in that leftover read at the right site, the quantities it already knows by
those names. The construction supplies the noun; physics supplies the reading; and the book keeps them apart on purpose, so
that no reading is ever mistaken for a derivation of the magnitude it names.

There is a unity here worth stating, because it is the reason a single instrument can speak of gravity and charge and phase
in one breath. All four readings are contractions or integrals of \emph{one} curvature, differing only in how the indices are
handled --- which are summed, which left free, which integrated over what. Mass is a trace, charge is a flux, phase is a loop,
curvature is the untouched tensor; and trace, flux, loop, and tensor are four operations on a single geometric object. This
is why the instrument does not need four separate machines. It needs one residue and four ways of reading it, and the ``four
faces'' the book will reach in Part III are precisely these four handlings of the one invariant. The economy is not a
convenience; it is the claim that mass, charge, and phase are not independent posits but one leftover seen from four sides.

In the two verbs, the residue is the funge that every later chapter tanges a reading out of. The curvature is bagged, whole
and invariant, the world's covariant leftover; and to read it as mass, or charge, or phase, is to tange one characteristic
out of that bag --- to select the trace, or the flux, or the loop --- while leaving the rest in place. The residue is the
common funge; the physical names are the tanges the instrument takes from it at different sites. So the antagonist that would
not disappear turns out to be the treasury the whole book draws on: one carried leftover, read at four places, becomes every
magnitude the instrument will later bracket. What remains is to walk the tower back down and find, at its seams, where the
residue is deposited --- and that descent is the next part.
